% !TeX program = pdflatex
\documentclass[12pt,a4paper]{article}

\newcommand{\notelanguage}{finnish}
% Shared preamble for course notes and exercise sets.

\usepackage[T1]{fontenc}
\usepackage{lmodern}
\usepackage[english]{babel}

\usepackage[a4paper,margin=30mm]{geometry}
\usepackage{microtype}
\usepackage{graphicx}
\usepackage{enumitem}

\usepackage{amsmath}
\usepackage{amssymb}
\usepackage{amsthm}
\usepackage{mathtools}
\usepackage{aliascnt}

\usepackage[hidelinks,pdfusetitle]{hyperref}

\numberwithin{equation}{section}
\setlist{itemsep=0.25em,topsep=0.5em}

\theoremstyle{plain}
\newtheorem{theorem}{Theorem}[section]
\newaliascnt{lemma}{theorem}
\newtheorem{lemma}[lemma]{Lemma}
\aliascntresetthe{lemma}
\newaliascnt{proposition}{theorem}
\newtheorem{proposition}[proposition]{Proposition}
\aliascntresetthe{proposition}
\newaliascnt{corollary}{theorem}
\newtheorem{corollary}[corollary]{Corollary}
\aliascntresetthe{corollary}

\theoremstyle{definition}
\newaliascnt{definition}{theorem}
\newtheorem{definition}[definition]{Definition}
\aliascntresetthe{definition}
\newaliascnt{example}{theorem}
\newtheorem{example}[example]{Example}
\aliascntresetthe{example}
\newaliascnt{problem}{theorem}
\newtheorem{problem}[problem]{Problem}
\aliascntresetthe{problem}

\theoremstyle{remark}
\newaliascnt{remark}{theorem}
\newtheorem{remark}[remark]{Remark}
\aliascntresetthe{remark}

\usepackage[nameinlink,noabbrev]{cleveref}

\crefname{theorem}{theorem}{theorems}
\crefname{lemma}{lemma}{lemmas}
\crefname{proposition}{proposition}{propositions}
\crefname{corollary}{corollary}{corollaries}
\crefname{definition}{definition}{definitions}
\crefname{example}{example}{examples}
\crefname{problem}{problem}{problems}
\crefname{remark}{remark}{remarks}

\DeclareMathOperator{\im}{im}
\DeclareMathOperator{\rank}{rank}
\DeclarePairedDelimiter{\abs}{\lvert}{\rvert}
\DeclarePairedDelimiter{\norm}{\lVert}{\rVert}
\DeclarePairedDelimiter{\set}{\lbrace}{\rbrace}


\title{Harjoituksen 2 ratkaisut}
\author{}
\date{}

\begin{document}

\maketitle

\noindent

\begin{exercise}
  Oletataan, että \(v(p_0) = 1\), \(v(p_1) = 0\) ja \(v(p_2) = 1\). Määritä
\begin{enumerate}[label=\alph*), nosep]
  \item \(v((p_0 \lor \neg p_1) \rightarrow (p_1 \land p_2))\),
  \item \(v((p_0 \leftrightarrow \neg p_1) \land \neg p_2)\).
\end{enumerate}
\end{exercise}

\begin{solution}
\leavevmode\par
\begin{enumerate}[label=\alph*), nosep]
\item Totuusarvon määritelmän mukaan,
  \[
    \begin{aligned}
      v((p_0 \lor \neg p_1) \rightarrow (p_1 \land p_2)) &= v((1 \lor
      \neg 0) \leftarrow (0 \land 1)) \\
      &= v((1 \lor 1) \leftarrow 0) \\
      &= v(1 \leftrightarrow 0) \\
      &= 0.
    \end{aligned}
  \]

\item Totuusarvon määritelmän mukaan,
  \[
    \begin{aligned}
      v((p_0 \leftrightarrow \neg p_1) \land \neg p_2) &= v((1
      \leftrightarrow \neg 0) \land \neg 1) \\
      &= v((1 \leftrightarrow 1) \land 0) \\
      &= v(1 \land 0) \\
      &= 0.
    \end{aligned}
  \]
\end{enumerate}
\end{solution}

\begin{exercise}
Anna esimerkki totuusjakaumasta, jolla lause
\begin{enumerate}[label=\alph*), nosep]
\item \(((p_0 \rightarrow p_1) \land \neg(p_1 \rightarrow p_0))\)
\item \(((p_0 \land p_1) \leftrightarrow (p_0 \lor p_1))\)
\end{enumerate}
on tosi.
\end{exercise}

\begin{solution}
\leavevmode\par
\begin{enumerate}[label=\alph*), nosep]
\item Olkoon \(v\) totuusjakauma, jolla \(v(p_0) = 0\) ja \(v(p_1) = 1\).
Laskemalla lauseen totuusarvo, saadaan
\[
\begin{aligned}
  v((p_0 \rightarrow p_1) \land \neg(p_1 \rightarrow p_0)) &=
  v((0 \rightarrow 1) \land \neg(1 \rightarrow 0)) \\
  &= v(1 \land 1) \\
  &= 1.
\end{aligned}
\]

\item Olkoon \(v\) totuusjakauma, jolla \(v(p_0) = v(p_1) = 0\).
Laskemalla lauseen totuusarvo, saadaan
\[
\begin{aligned}
  v((p_0 \land p_1) \leftrightarrow (p_0 \lor p_1)) &=
  v((0 \land 0) \leftrightarrow (0 \lor 0)) \\
  &= v(0 \leftrightarrow 0) \\
  &= 1.
\end{aligned}
\]
\end{enumerate}
\end{solution}

\begin{exercise}
\end{exercise}

\begin{solution}
\end{solution}

\begin{exercise}
\end{exercise}

\begin{solution}
\end{solution}

\begin{exercise}
\end{exercise}

\begin{solution}
\end{solution}

\begin{exercise}
\end{exercise}

\begin{solution}
\end{solution}

\begin{exercise}
\end{exercise}

\begin{solution}
\end{solution}

\begin{exercise}
\end{exercise}

\begin{solution}
\end{solution}

\begin{exercise}
\end{exercise}

\begin{solution}
\end{solution}

\begin{exercise}
\end{exercise}

\begin{solution}
\end{solution}

\begin{exercise}[Haastetehtävä]
\end{exercise}

\begin{solution}
\end{solution}

\end{document}
